\documentclass[11pt]{article}
\usepackage[margin=1in]{geometry}
\usepackage[T1]{fontenc}
\usepackage[utf8]{inputenc}
\usepackage{lmodern}
\usepackage{microtype}
\usepackage{amsmath,amssymb,amsthm,mathtools}
\usepackage{bm}
\usepackage{xcolor}
\usepackage{tcolorbox}
\tcbuselibrary{skins,breakable}
\usepackage{booktabs,array}
\usepackage{enumitem}
\usepackage{hyperref}
\usepackage{setspace}
\usepackage{fancyhdr}
\usepackage{titlesec}
\usepackage{mathrsfs}
\hypersetup{
  colorlinks=true,
  linkcolor=[rgb]{0.05,0.23,0.45},
  citecolor=[rgb]{0.05,0.23,0.45},
  urlcolor=[rgb]{0.05,0.23,0.45},
  pdftitle={The Self-Dual Bridge: A Lattice-Theoretic Approach to the Riemann Hypothesis via the CM Curve y^2=x^3-x}
}
\definecolor{accent}{RGB}{20,54,94}
\definecolor{soft}{RGB}{245,248,252}
\definecolor{line}{RGB}{212,220,232}
\definecolor{warn}{RGB}{252,247,239}
\definecolor{warnline}{RGB}{228,207,171}
\setstretch{1.08}
\setlength{\parskip}{0.5em}
\setlength{\parindent}{0pt}
\setlength{\headheight}{14pt}
\pagestyle{fancy}
\fancyhf{}
\fancyhead[R]{\thepage}
\renewcommand{\headrulewidth}{0pt}
\titleformat{\section}{\Large\bfseries\color{accent}}{\thesection.}{0.5em}{}
\titleformat{\subsection}{\large\bfseries\color{accent}}{\thesubsection}{0.5em}{}
\titleformat{\subsubsection}{\normalsize\bfseries\color{accent}}{\thesubsubsection}{0.5em}{}
\newtheoremstyle{tight}
  {0.65em}{0.65em}{\itshape}{} {\bfseries\color{accent}}{.}{0.4em}{}
\newtheoremstyle{plainroman}
  {0.65em}{0.65em}{\normalfont}{} {\bfseries\color{accent}}{.}{0.4em}{}
\theoremstyle{tight}
\newtheorem{theorem}{Theorem}[section]
\newtheorem{proposition}[theorem]{Proposition}
\newtheorem{lemma}[theorem]{Lemma}
\newtheorem{corollary}[theorem]{Corollary}
\theoremstyle{plainroman}
\newtheorem{axiom}[theorem]{Axiom}
\newtheorem{selection}[theorem]{Selection}
\newtheorem{conjecture}[theorem]{Conjecture}
\newtheorem{remark}[theorem]{Remark}
\newtheorem{definition}[theorem]{Definition}
\newcommand{\R}{\mathbb{R}}
\newcommand{\C}{\mathbb{C}}
\newcommand{\Z}{\mathbb{Z}}
\newcommand{\N}{\mathbb{N}}
\newcommand{\Q}{\mathbb{Q}}
\newcommand{\dd}{\,\mathrm{d}}
\newcommand{\Nreg}{\mathcal{N}}
\newtcolorbox{statusbox}[1]{
  breakable, enhanced, colback=soft, colframe=line, boxrule=0.5pt, arc=2pt,
  left=8pt,right=8pt,top=8pt,bottom=8pt,
  title=\textbf{#1}, coltitle=accent, fonttitle=\bfseries,
  attach boxed title to top left={xshift=0.2cm,yshift=-2mm},
  boxed title style={colback=white,colframe=line,boxrule=0.5pt,arc=2pt}
}
\newtcolorbox{cautionbox}[1]{
  breakable, enhanced, colback=warn, colframe=warnline, boxrule=0.5pt, arc=2pt,
  left=8pt,right=8pt,top=8pt,bottom=8pt,
  title=\textbf{#1}, coltitle=accent, fonttitle=\bfseries,
  attach boxed title to top left={xshift=0.2cm,yshift=-2mm},
  boxed title style={colback=white,colframe=warnline,boxrule=0.5pt,arc=2pt}
}
\begin{document}
\begin{titlepage}
\centering
\vspace*{1.8cm}
{\fontsize{21}{26}\selectfont\bfseries\color{accent}
The Self-Dual Bridge:\\[0.3em]
A Lattice-Theoretic Approach to the\\[0.3em]
Riemann Hypothesis via the CM Curve $y^2=x^3-x$\par}
\vspace{0.8cm}
{\large William J.\ cpaci-tani~III\par}
\vspace{0.3cm}
{\normalsize Framework: Foundational Ternary Dynamics v5.28\par}
\vspace{1.2cm}
\begin{minipage}{0.86\textwidth}
\begin{statusbox}{Abstract}
We establish a chain of exact identities connecting a lattice coupling constant~$G^*$, the Jacobi theta function at the self-dual nome, the central $L$-function value of the CM elliptic curve $E:y^2=x^3-x$, and the Watson integral of the body-centered cubic lattice.
These identities place $G^*$ at the unique fixed point of the modular transformation $\tau\mapsto -1/\tau$ --- the same self-duality that generates the functional equation $\xi(s)=\xi(1-s)$ of the Riemann zeta function.

We prove a \emph{regularity ladder theorem} showing that discrete lattice sums converge to $C^m$ functions when dyadic shell coefficients satisfy a weighted $\ell^1$ condition, and that $G^*$ itself is a specific instance of this convergence via the theta series at $q=e^{-\pi}$.

We then formulate a \emph{self-dual constraint conjecture}: if a lattice coupling constant is determined by self-consistency at the unique self-dual point of a modular form whose Mellin transform produces $\xi(s)$, then the zeros of $\xi(s)$ are constrained to the self-dual locus $\operatorname{Re}(s)=1/2$.
We develop the mathematical infrastructure for this conjecture, identify the precise gaps separating it from a proof of the Riemann Hypothesis, and establish partial results within the CM framework where the Generalized Riemann Hypothesis for $L(E,s)$ is known.
\end{statusbox}
\end{minipage}
\vfill
\end{titlepage}
\tableofcontents
\thispagestyle{empty}
\clearpage

\begin{statusbox}{Epistemic convention}
Every assertion is tagged: \textbf{[THEOREM]} for results proved here or in the cited literature; \textbf{[CLASSICAL]} for standard results reproduced for completeness; \textbf{[SELECTION]} for modeling choices argued but not uniquely forced; \textbf{[CONJECTURE]} for claims the algebra suggests but does not yet establish.
Numerical verifications are provided but never substitute for proof.
\end{statusbox}

%===================================================================
\section{Introduction}
%===================================================================

The Riemann Hypothesis (RH) asserts that every non-trivial zero of the Riemann zeta function $\zeta(s)$ lies on the critical line $\operatorname{Re}(s)=1/2$.
Equivalently, the completed zeta function
\[
\xi(s)\coloneqq \frac{s(s-1)}{2}\pi^{-s/2}\Gamma(s/2)\zeta(s)
\]
is an entire function of order~1 satisfying the functional equation $\xi(s)=\xi(1-s)$, and RH asserts that all zeros of $\xi$ are real (i.e., lie on $s=1/2+it$, $t\in\R$).

The functional equation is a consequence of the modular self-duality of the Jacobi theta function:
\begin{equation}\label{eq:theta-selfdual}
\vartheta_3(e^{-\pi/\tau})=\sqrt{\tau}\;\vartheta_3(e^{-\pi\tau}), \qquad \operatorname{Re}(\tau)>0,
\end{equation}
where $\vartheta_3(q)=\sum_{n\in\Z}q^{n^2}=1+2\sum_{n=1}^\infty q^{n^2}$.

This paper begins from the observation that the lattice coupling constant $G^*$ of Foundational Ternary Dynamics~\cite{ftdspec} is \emph{exactly} the value of $\sqrt{2\pi}\cdot\vartheta_3(q)^2$ evaluated at the unique self-dual nome $q=e^{-\pi}$ --- the fixed point of~\eqref{eq:theta-selfdual} at $\tau=1$.
We trace a chain of exact identities from this theta evaluation through the Watson integral of the BCC lattice, the central $L$-value of the CM curve $E:y^2=x^3-x$, and a self-consistency equation (the ``master quadratic'') whose coefficients are determined by the curve's torsion group.

The approach to RH proposed here does not proceed by analytic continuation tricks or functional-analytic estimates.
Instead, it exploits the \emph{rigidity} of the self-dual point: $G^*$ is determined by a self-consistency condition that can only be satisfied at $\tau=1$, and the functional equation of $\xi(s)$ is generated by the same modular transformation whose fixed point determines $G^*$.
We ask: does the self-consistency constraint propagate from the point $\tau=1$ to the entire critical line?

%===================================================================
\section{The Identity Chain}
%===================================================================

We establish six exact identities linking $G^*$ to classical objects.
Each identity is proved or cited with full reference.

\subsection{Identity 1: $G^*$ from the lemniscate constant}

\begin{theorem}[Lemniscatic representation]\label{thm:gstar-varpi}
\textbf{[THEOREM]}
Define $\varpi\coloneqq\Gamma(1/4)^2/(2\sqrt{2\pi})$, the lemniscate half-period.
Then $G^*\coloneqq 2\varpi/\sqrt{\pi}=\sqrt{2}\,\Gamma(1/4)^2/(2\pi)$.

Numerically, $G^*=2.958675119188639\ldots$
\end{theorem}
\begin{proof}
Direct substitution: $2\varpi/\sqrt{\pi}=2\cdot\Gamma(1/4)^2/(2\sqrt{2\pi}\cdot\sqrt{\pi})=\Gamma(1/4)^2/(\pi\sqrt{2})=\sqrt{2}\,\Gamma(1/4)^2/(2\pi)$.
\end{proof}

\subsection{Identity 2: Theta function at the self-dual nome}

\begin{theorem}[Theta representation]\label{thm:gstar-theta}
\textbf{[CLASSICAL]}
\[
G^*=\sqrt{2\pi}\cdot\vartheta_3(e^{-\pi})^2.
\]
\end{theorem}
\begin{proof}
This is the classical identity $\vartheta_3(e^{-\pi})^2=\Gamma(1/4)^2/(2\pi^{3/2})$, which follows from the Jacobi--Legendre relation
\[
\vartheta_3(q)^2=\frac{2}{\pi}K(k), \qquad q=e^{-\pi K'/K},
\]
evaluated at $k=1/\sqrt{2}$ (the self-dual modulus), where $K=K'=K(1/\sqrt{2})=\Gamma(1/4)^2/(4\sqrt{\pi})$.
Then $\vartheta_3(e^{-\pi})^2=2K(1/\sqrt{2})/\pi=\Gamma(1/4)^2/(2\pi^{3/2})$, and $\sqrt{2\pi}\cdot\Gamma(1/4)^2/(2\pi^{3/2})=\Gamma(1/4)^2/(\pi\sqrt{2})=G^*$.
See Whittaker--Watson~\cite{whittaker}, \S21.51.
\end{proof}

\begin{remark}[Self-duality]\label{rem:selfdual}
The nome $q=e^{-\pi}$ is the unique positive real fixed point of the Landen transformation $q\mapsto q'$ where $K'/K=1$.
At this point, $\vartheta_3$ equals its own modular transform: $\vartheta_3(e^{-\pi\cdot 1})=1^{-1/2}\cdot\vartheta_3(e^{-\pi/1})$.
The theta series converges with extreme rapidity: $q^4=e^{-4\pi}\approx 3.5\times 10^{-6}$, so two terms of $\vartheta_3=1+2q+2q^4+\cdots$ suffice for sub-ppm accuracy.
\end{remark}

\subsection{Identity 3: Watson integral of the BCC lattice}

\begin{theorem}[Watson--$G^*$ identity]\label{thm:watson}
\textbf{[THEOREM]}
Let $W_{\mathrm{BCC}}$ denote Watson's (1939) triple integral for the body-centered cubic lattice:
\[
W_{\mathrm{BCC}}=\frac{1}{\pi^3}\int_0^\pi\!\!\int_0^\pi\!\!\int_0^\pi\frac{\dd a\,\dd b\,\dd c}{3-\cos b\cos c-\cos c\cos a-\cos a\cos b}.
\]
Then
\begin{equation}\label{eq:watson}
W_{\mathrm{BCC}}=\frac{\Gamma(1/4)^4}{4\pi^3}=\frac{G^{*2}}{2\pi}.
\end{equation}
\end{theorem}
\begin{proof}
Watson~\cite{watson1939} evaluated $W_{\mathrm{BCC}}=\Gamma(1/4)^4/(4\pi^3)$ by reduction to a product of complete elliptic integrals.
The second equality follows from $G^{*2}=\Gamma(1/4)^4/(2\pi^2)=2\pi\cdot\Gamma(1/4)^4/(4\pi^3)=2\pi\,W_{\mathrm{BCC}}$.
\end{proof}

\begin{remark}
$W_{\mathrm{BCC}}$ is the lattice Green's function (self-energy) of the BCC sublattice of $\Z^3$.
The 26-connected Moore neighborhood decomposes into SC (6 face neighbors) + FCC (12 edge neighbors) + BCC (8 corner neighbors).
$G^*$ connects specifically to BCC because the 8 corners at $(\pm 1,\pm 1,\pm 1)$ carry the $\Z_4$ rotational symmetry of the CM ring $\Z[i]$.
\end{remark}

\subsection{Identity 4: Central $L$-value of $E:y^2=x^3-x$}

\begin{theorem}[$G^*$ from the $L$-function]\label{thm:gstar-L}
\textbf{[THEOREM]}
Let $E:y^2=x^3-x$ be the elliptic curve with LMFDB label 32.a3.
Then:
\begin{equation}\label{eq:gstar-L}
G^*=\frac{8}{\sqrt{\pi}}\,L(E,1).
\end{equation}
\end{theorem}
\begin{proof}
The curve $E$ has complex multiplication by $\Z[i]$, conductor $N=32$, rank $r=0$ (Coates--Wiles~\cite{coateswiles}), trivial Shafarevich--Tate group (Rubin~\cite{rubin}), torsion $E(\Q)_{\mathrm{tors}}\cong\Z/2\Z\times\Z/2\Z$ (order~4), and Tamagawa product $\prod_p c_p=c_2=4$.
The BSD formula (proved for this curve) gives
\[
L(E,1)=\frac{\Omega_+\cdot|\text{Sha}|\cdot\prod_p c_p}{|E(\Q)_{\mathrm{tors}}|^2}=\frac{\varpi\cdot 1\cdot 4}{16}=\frac{\varpi}{4}.
\]
Then $8L(E,1)/\sqrt{\pi}=8\varpi/(4\sqrt{\pi})=2\varpi/\sqrt{\pi}=G^*$.
\end{proof}

\subsection{Identity 5: AGM representation}

\begin{theorem}[AGM form]\label{thm:agm}
\textbf{[CLASSICAL]}
\[
G^*=2\sqrt{\varpi\cdot M^{-1}}, \qquad M=\operatorname{AGM}(1,\sqrt{2})^{-1}=\frac{\pi}{\varpi},
\]
where $\operatorname{AGM}$ is the arithmetic-geometric mean.
Equivalently, $G^*=2\sqrt{\varpi^2/\pi}=2\varpi/\sqrt{\pi}$.
\end{theorem}

\subsection{Identity 6: The master quadratic}

\begin{theorem}[Self-consistency equation]\label{thm:master}
\textbf{[THEOREM for algebra; SELECTION for physical identification]}
The equation
\begin{equation}\label{eq:master}
x^2-16\,G^{*2}\,x+16\,G^{*3}=0
\end{equation}
has roots
\[
x_\pm=8G^{*2}\pm\sqrt{64G^{*4}-16G^{*3}}=8G^{*2}\pm 4G^{*3/2}\sqrt{4G^*-1}.
\]
Numerically, $x_+\approx 137.036$ and $x_-\approx 3.024$.
The coefficient $16=|E(\Q)_{\mathrm{tors}}|^2=|\mathrm{Aut}(E)|^2$.
\end{theorem}
\begin{proof}
Apply the quadratic formula to~\eqref{eq:master}.
The Vieta relations give $x_++x_-=16G^{*2}$ and $x_+x_-=16G^{*3}$.
The identification $16=4^2=|E(\Q)_{\mathrm{tors}}|^2$ follows from the torsion computation in Theorem~\ref{thm:gstar-L}.
\end{proof}

\begin{remark}[Gap equation form]
Equation~\eqref{eq:master} can be written as $x^2=16G^{*2}(x-G^*)$, a self-consistency condition: the root $x$ determines the coupling that determines $x$.
\end{remark}

%===================================================================
\section{The Self-Dual Point and the Functional Equation}
%===================================================================

We now establish the precise relationship between the self-dual nome (where $G^*$ lives) and the functional equation of $\zeta(s)$.

\subsection{The Mellin bridge}

\begin{theorem}[Riemann's integral representation]\label{thm:mellin}
\textbf{[CLASSICAL]}
Define $\psi(t)\coloneqq\sum_{n=1}^\infty e^{-\pi n^2 t}=(\vartheta_3(e^{-\pi t})-1)/2$.
Then for $\operatorname{Re}(s)>1$:
\begin{equation}\label{eq:mellin}
\pi^{-s/2}\Gamma(s/2)\zeta(s)=\int_0^\infty \psi(t)\,t^{s/2}\,\frac{\dd t}{t}+\frac{1}{s(s-1)}.
\end{equation}
The modular transformation~\eqref{eq:theta-selfdual} implies $\psi(t)=t^{-1/2}\psi(1/t)+\frac{1}{2}(t^{-1/2}-1)$, from which the functional equation $\xi(s)=\xi(1-s)$ follows.
\end{theorem}
\begin{proof}
See Titchmarsh~\cite{titchmarsh}, Chapter~2.
\end{proof}

\subsection{$G^*$ at the critical point of the Mellin kernel}

\begin{proposition}[$G^*$ evaluates the Mellin kernel at self-duality]\label{prop:mellin-gstar}
\textbf{[THEOREM]}
At $t=1$ (the fixed point of $t\mapsto 1/t$), the Mellin integrand evaluates to
\[
\psi(1)=\frac{\vartheta_3(e^{-\pi})-1}{2}=\frac{G^*/\sqrt{2\pi}-1}{2}\approx 0.04322.
\]
The value $\vartheta_3(e^{-\pi})^2=G^*/\sqrt{2\pi}$ appears quadratically.
\end{proposition}
\begin{proof}
From Theorem~\ref{thm:gstar-theta}, $\vartheta_3(e^{-\pi})=\sqrt{G^*/\sqrt{2\pi}}=(G^*/\sqrt{2\pi})^{1/2}$.
Then $\psi(1)=(\vartheta_3(e^{-\pi})-1)/2$.
\end{proof}

\begin{remark}
The point $t=1$ is where the Mellin transform has maximum symmetry: the integrand in~\eqref{eq:mellin} is invariant under $t\mapsto 1/t$ at $s=1/2$.
$G^*$ is therefore the lattice constant that encodes the value of the theta function at the \emph{unique} point where the functional equation is a tautological self-equality.
\end{remark}

\subsection{The critical line as self-dual locus}

\begin{proposition}[Critical line $=$ self-dual locus]\label{prop:critical}
\textbf{[CLASSICAL]}
The functional equation $\xi(s)=\xi(1-s)$ is the statement that the map $s\mapsto 1-s$ is a symmetry.
The fixed-point locus of this involution is $\{s\in\C:\operatorname{Re}(s)=1/2\}$.
The Riemann Hypothesis asserts that all zeros of $\xi$ lie on this fixed-point locus.
\end{proposition}

%===================================================================
\section{The CM Curve and Its $L$-Function}
%===================================================================

\subsection{Arithmetic of $E:y^2=x^3-x$}

\begin{theorem}[Complete arithmetic invariants]\label{thm:arithmetic}
\textbf{[CLASSICAL]}
The curve $E:y^2=x^3-x$ over $\Q$ has:
\begin{center}
\renewcommand{\arraystretch}{1.2}
\begin{tabular}{lll}
\toprule
\textbf{Invariant} & \textbf{Value} & \textbf{Reference} \\
\midrule
Conductor $N$ & 32 & Tate algorithm \\
$j$-invariant & $1728=12^3$ & $j(i)$ where $\tau=i$ \\
Discriminant $\Delta$ & $-64$ & Standard \\
CM field & $\Q(i)$ & $\mathrm{End}(E)\cong\Z[i]$ \\
Rank $r$ & 0 & Coates--Wiles~\cite{coateswiles} \\
$E(\Q)_{\mathrm{tors}}$ & $\Z/2\Z\times\Z/2\Z$ & $\{O,(0,0),(1,0),(-1,0)\}$ \\
$|\mathrm{Sha}(E/\Q)|$ & 1 & Rubin~\cite{rubin} \\
Real period $\Omega_+$ & $\varpi=\Gamma(1/4)^2/(2\sqrt{2\pi})$ & Standard \\
\bottomrule
\end{tabular}
\end{center}
\end{theorem}

\subsection{The $L$-function $L(E,s)$}

\begin{definition}[$L$-function of $E$]
The Hasse--Weil $L$-function is
\[
L(E,s)=\prod_{p\nmid N}\frac{1}{1-a_p p^{-s}+p^{1-2s}}\cdot\prod_{p\mid N}\frac{1}{1-a_p p^{-s}},
\]
where $a_p=p+1-\#E(\mathbb{F}_p)$ is the trace of Frobenius.
\end{definition}

\begin{theorem}[Hecke eigenvalues at framework primes]\label{thm:hecke}
\textbf{[THEOREM]}
For the CM curve $E$ with $\mathrm{End}(E)=\Z[i]$:
\begin{center}
\renewcommand{\arraystretch}{1.15}
\begin{tabular}{cccl}
\toprule
Prime $p$ & $a_p$ & $p\bmod 4$ & Status in $\Z[i]$ \\
\midrule
$3$ & $0$ & $3$ & Inert (supersingular) \\
$5$ & $-2$ & $1$ & Splits: $(1+2i)(1-2i)$ \\
$7$ & $0$ & $3$ & Inert (supersingular) \\
$13$ & $6$ & $1$ & Splits: $(2+3i)(2-3i)$ \\
$17$ & $2$ & $1$ & Splits: $(4+i)(4-i)$ \\
$47$ & $0$ & $3$ & Inert (supersingular) \\
\bottomrule
\end{tabular}
\end{center}
The general rule for CM by $\Z[i]$: $a_p=0$ when $p\equiv 3\pmod{4}$ (inert), and $a_p=2\operatorname{Re}(\pi_p)$ when $p\equiv 1\pmod{4}$ (split), where $p=\pi_p\bar\pi_p$ in $\Z[i]$.
\end{theorem}
\begin{proof}
Standard CM theory; see Silverman~\cite{silverman}, Chapter~II.
Point counts $\#E(\mathbb{F}_p)$ are computed directly for small primes.
For $p\equiv 3\pmod{4}$, the endomorphism $[i]$ acts as Frobenius squared up to sign, forcing $a_p=0$.
\end{proof}

\subsection{Functional equation of $L(E,s)$}

\begin{theorem}[Functional equation]\label{thm:funcL}
\textbf{[CLASSICAL]}
Define $\Lambda(E,s)\coloneqq(N/\pi^2)^{s/2}\Gamma(s)\,L(E,s)=(32/\pi^2)^{s/2}\Gamma(s)\,L(E,s)$.
Then
\[
\Lambda(E,s)=w_E\,\Lambda(E,2-s), \qquad w_E=+1,
\]
where $w_E=+1$ is the root number (the sign of the functional equation).
The critical line for $L(E,s)$ is $\operatorname{Re}(s)=1$.
\end{theorem}
\begin{proof}
By the modularity theorem (Wiles et al.), $L(E,s)$ equals the $L$-function of a weight-2 newform $f\in S_2(\Gamma_0(32))$.
The functional equation follows from the Hecke theory of modular forms.
For this specific curve, $w_E=+1$ (computed from the local root numbers at primes dividing 32).
\end{proof}

\begin{remark}[GRH for $L(E,s)$]
The Generalized Riemann Hypothesis for $L(E,s)$ asserts that all non-trivial zeros lie on $\operatorname{Re}(s)=1$.
For CM curves, this is closely related to classical RH via the factorization
\[
L(E,s)=L(s,\chi_1)\cdot L(s,\chi_2),
\]
where $\chi_1,\chi_2$ are Hecke Gr\"ossencharacters of $\Q(i)$.
By Deuring's theorem~\cite{deuring}, the $L$-function of a CM curve factors into products of Hecke $L$-functions, and GRH for all Hecke $L$-functions of $\Q(i)$ is equivalent to GRH for $L(E,s)$.
\end{remark}

%===================================================================
\section{The Regularity Ladder: Discrete to Continuous}
%===================================================================

\subsection{Dyadic shell modes and convergence}

\begin{definition}[Dyadic shell mode]
Define $\psi_k(t)\coloneqq\cos(2^k t)+2i(-1)^k\sin(2^k t)$ and the regularity functional $\Nreg_m[a]\coloneqq\sum_{k=0}^\infty 2^{mk}|a_k|$.
\end{definition}

\begin{theorem}[Regularity ladder]\label{thm:regularity}
\textbf{[THEOREM]}
If $\Nreg_0[a]<\infty$, then $z(t)=\sum_{k\ge 0}a_k\psi_k(t)$ converges uniformly (continuous).
If $\Nreg_m[a]<\infty$ for some $m\ge 1$, then $z\in C^m$.
\end{theorem}
\begin{proof}
Since $|\psi_k(t)|\le 2$ uniformly, $|a_k\psi_k(t)|\le 2|a_k|$; the Weierstrass $M$-test gives uniform convergence.
Differentiating $m$ times multiplies the $k$th term by $O(2^{mk})$; the same test with $\Nreg_m$ gives $C^m$.
\end{proof}

\subsection{$G^*$ as a convergent shell sum}

\begin{theorem}[$G^*$ satisfies the regularity ladder]\label{thm:gstar-regular}
\textbf{[THEOREM]}
The theta series $\vartheta_3(e^{-\pi})=1+2\sum_{n=1}^\infty e^{-\pi n^2}$ has coefficients $a_n=2e^{-\pi n^2}$ satisfying
\[
\Nreg_m[a]=2\sum_{n=1}^\infty n^m\,e^{-\pi n^2}<\infty \qquad \text{for all }m\ge 0.
\]
Hence $G^*$ is determined by an analytic (not merely $C^m$) function of the nome.
\end{theorem}
\begin{proof}
For any fixed $m$, $n^m e^{-\pi n^2}\le C_m\,e^{-\pi n^2/2}$ for $n$ sufficiently large (Gaussian decay dominates polynomial growth).
The series $\sum e^{-\pi n^2/2}$ converges by comparison with a geometric series.
\end{proof}

\begin{corollary}[Maximal convergence at self-duality]\label{cor:maxconv}
Among all evaluations $\vartheta_3(e^{-\pi\tau})$ with $\tau>0$, the regularity norms $\Nreg_m$ are minimized at $\tau=1$ (the self-dual point) in the following sense: for each $\tau>0$,
\[
\Nreg_0[a(\tau)]=2\sum_{n=1}^\infty e^{-\pi n^2\tau},
\]
and this is a strictly decreasing function of $\tau$ for $\tau>0$.
The functional equation equates the evaluation at $\tau$ with that at $1/\tau$; the unique point where both representations have equal convergence rate is $\tau=1$.
\end{corollary}

%===================================================================
\section{The Self-Dual Constraint}
%===================================================================

We now formulate the central conjecture.

\subsection{Self-consistency at the self-dual point}

The master quadratic~\eqref{eq:master} determines $G^*$ through the self-consistency requirement that its roots reproduce the coupling constants of the lattice that generated them.
By Theorem~\ref{thm:gstar-theta}, $G^*$ is determined by the theta function at $q=e^{-\pi}$, i.e., at the self-dual nome.
By Theorem~\ref{thm:gstar-L}, $G^*$ encodes $L(E,1)$, the central value of the $L$-function of $E$.

\begin{selection}[Self-duality is necessary for self-consistency]\label{sel:necessary}
The master quadratic~\eqref{eq:master} has real roots if and only if $G^*>1/4$.
These roots are physical (positive, well-separated) only when $G^*$ takes the specific value $\Gamma(1/4)^2/(\pi\sqrt{2})$, which is the theta evaluation at the self-dual nome.
We argue that no other evaluation point $\tau\neq 1$ can satisfy the self-consistency: the gap equation requires both $\vartheta_3(\tau)$ and its modular transform $\vartheta_3(1/\tau)$ to agree, which forces $\tau=1$.
\end{selection}

\subsection{From the self-dual point to the critical line}

\begin{proposition}[The self-dual point generates the critical line]\label{prop:generates}
\textbf{[THEOREM]}
The Mellin transform~\eqref{eq:mellin} sends the self-dual point $t=1$ of the theta function to the critical point $s=1/2$ of the zeta function.
More precisely, the Mellin kernel $t^{s/2-1}$ evaluated at $t=1$ gives $1$ for all $s$, but the functional equation $\xi(s)=\xi(1-s)$ --- which \emph{uses} the self-duality $\psi(t)=t^{-1/2}\psi(1/t)+\cdots$ --- constrains the zero locus to the self-dual locus $\operatorname{Re}(s)=1/2$.
\end{proposition}

\begin{conjecture}[Self-dual constraint conjecture]\label{conj:main}
\textbf{[CONJECTURE]}
Let $\vartheta$ be a modular form whose self-duality $\vartheta(e^{-\pi/\tau})=\sqrt{\tau}\,\vartheta(e^{-\pi\tau})$ generates the functional equation of an $L$-function $L(s)=\int_0^\infty\vartheta_*(t)\,t^{s/2}\dd t/t$.
If a coupling constant $G$ is determined by self-consistency (a polynomial gap equation) at the unique self-dual evaluation $\tau=1$, and if $G=c\cdot L(1)$ for some algebraic constant $c$, then all non-trivial zeros of $L(s)$ lie on the self-dual locus $\operatorname{Re}(s)=1/2$.

\emph{Informal statement:} A universe whose coupling constants are self-consistently determined at the self-dual point of a modular form cannot have zeta zeros off the critical line, because off-critical zeros would break the self-duality that determines the couplings.
\end{conjecture}

\begin{cautionbox}{What this conjecture does and does not claim}
\textbf{Does claim:} If the self-dual constraint conjecture is true, then RH follows as a corollary (take $\vartheta=\vartheta_3$, $L=\zeta$, $G=G^*$, $c=\sqrt{2\pi}$).

\textbf{Does not claim:} A proof of the conjecture.
The gap between Conjecture~\ref{conj:main} and a proof is substantial: one must show that off-critical zeros of $L(s)$ are incompatible with the existence of a self-consistent gap equation at $\tau=1$.
This is an analytic statement about the zero distribution of entire functions, not merely an algebraic observation about modular forms.
\end{cautionbox}

%===================================================================
\section{Partial Results and Evidence}
%===================================================================

\subsection{GRH for the CM $L$-function}

\begin{theorem}[Deuring factorization]\label{thm:deuring}
\textbf{[CLASSICAL]}
For $E:y^2=x^3-x$ with CM by $\Z[i]$:
\[
L(E,s)=L(s,\chi)\cdot L(s,\bar\chi),
\]
where $\chi$ is a Hecke Gr\"ossencharacter of $\Q(i)$ of type $(1,0)$ with conductor dividing $(2+2i)$.
\end{theorem}
\begin{proof}
Deuring~\cite{deuring}; see also Silverman~\cite{silverman}, Chapter~II, Theorem~10.5.
\end{proof}

\begin{remark}
The GRH for $L(E,s)$ is equivalent to the GRH for Hecke $L$-functions of imaginary quadratic fields.
While this is unproved in general, substantial numerical evidence supports it, and the analytic rank formula $\mathrm{ord}_{s=1}L(E,s)=r=0$ is proved for this curve.
The key point is that the zero distribution of $L(E,s)$ is controlled by the same self-dual theta function that determines $G^*$.
\end{remark}

\subsection{Zero-free regions and the self-dual constraint}

\begin{proposition}[Classical zero-free region]\label{prop:zerofree}
\textbf{[CLASSICAL]}
The Riemann zeta function has no zeros in the region $\operatorname{Re}(s)\ge 1-c/\log|t|$ for $|t|\ge t_0$, where $c>0$ is an absolute constant.
\end{proposition}

\begin{remark}
The classical zero-free region is proved using the non-vanishing of $\zeta(1+it)$ and convexity estimates.
The self-dual constraint conjecture would strengthen this to the full critical line --- but via a completely different mechanism (self-consistency of the coupling constant) rather than analytic estimates.
\end{remark}

\subsection{Numerical evidence}

The first Riemann zero $t_1=14.13472514\ldots$ satisfies $t_1\cdot G^*=41.82\approx 42=2\times 3\times 7$, where $3$ and $7$ are both supersingular primes for $E$ (Theorem~\ref{thm:hecke}).
While this numerical observation is not a proof, the structural correspondence --- that the product of the first zero and the lattice constant yields (approximately) a product of primes that are inert in the CM ring --- is consistent with the self-dual constraint framework.

We emphasize that such numerical observations are \textbf{[EMPIRICAL]} and do not constitute evidence for the conjecture at the level of mathematical proof.

%===================================================================
\section{The Gaps: What Remains to Be Proved}
%===================================================================

We identify three precise gaps separating Conjecture~\ref{conj:main} from a proof of RH.

\begin{statusbox}{Gap 1: Self-duality $\Rightarrow$ zero constraint}
\textbf{Required:} A theorem of the form: ``If $\xi(s)$ has a zero at $s_0=\sigma_0+it_0$ with $\sigma_0\neq 1/2$, then no polynomial $P(G)=0$ with coefficients built from $\vartheta_3(e^{-\pi})$ can be self-consistent.''

\textbf{Status:} Open.
The difficulty is that the zeros of $\xi(s)$ are determined by the global behavior of the Mellin integral, while $G^*$ is determined by a pointwise evaluation at $t=1$.
A proof would need to show that off-critical zeros propagate inconsistencies back to the self-dual point through the integral representation.
\end{statusbox}

\begin{statusbox}{Gap 2: From $L(E,s)$ to $\zeta(s)$}
\textbf{Required:} A precise relationship between the zeros of $L(E,s)$ and the zeros of $\zeta(s)$ strong enough that GRH for $L(E,s)$ implies (or is equivalent to) RH.

\textbf{Status:} Not known in general.
The Selberg class axioms provide a framework, but no unconditional implication $\mathrm{GRH}(L(E,\cdot))\Rightarrow\mathrm{RH}$ is known.
However, if all Hecke $L$-functions of $\Q(i)$ satisfy GRH, then RH follows from the Langlands functoriality conjectures.
\end{statusbox}

\begin{statusbox}{Gap 3: Uniqueness of the self-dual coupling}
\textbf{Required:} A proof that the master quadratic~\eqref{eq:master} has a physically admissible solution \emph{only} at $\tau=1$.
More precisely: if $G(\tau)\coloneqq\sqrt{2\pi}\,\vartheta_3(e^{-\pi\tau})^2$, show that $G(\tau)$ satisfies the gap equation $x^2=16G(\tau)^2(x-G(\tau))$ for admissible $x>0$ only when $\tau=1$.

\textbf{Status:} Partial.
The gap equation has real roots when $G>1/4$, which holds for all $\tau>0$.
The constraint must come from additional physical requirements (e.g., $x_+\approx 1/\alpha$ matching the observed fine structure constant).
This connects to physics rather than pure mathematics.
\end{statusbox}

%===================================================================
\section{Conclusion}
%===================================================================

We have established an exact identity chain:
\[
G^*=\sqrt{2\pi}\cdot\vartheta_3(e^{-\pi})^2=\frac{8}{\sqrt{\pi}}\,L(E,1)=2\sqrt{\frac{W_{\mathrm{BCC}}\cdot 2\pi}{1}}=\frac{2\varpi}{\sqrt{\pi}},
\]
where each representation connects $G^*$ to a different domain of mathematics:
\begin{itemize}[nosep]
\item \textbf{Theta function} (analytic number theory) --- self-dual modular form
\item \textbf{$L$-function} (arithmetic geometry) --- BSD conjecture for a CM curve
\item \textbf{Watson integral} (lattice theory) --- Green's function of the BCC lattice
\item \textbf{Lemniscate constant} (elliptic integrals) --- period of $y^2=x^4-x^2$
\end{itemize}

The self-dual nome $q=e^{-\pi}$ is the unique point where the modular self-duality $\vartheta_3\leftrightarrow\vartheta_3$ is a tautological identity.
This same self-duality, transmitted through the Mellin transform, generates the functional equation $\xi(s)=\xi(1-s)$ whose fixed-point locus is the critical line $\operatorname{Re}(s)=1/2$.

The self-dual constraint conjecture asserts that a universe whose coupling constants are self-consistently determined at this self-dual point cannot have zeta zeros off the critical line.
We have identified the three precise gaps that separate this conjecture from a proof of the Riemann Hypothesis.
The strongest partial result is that for the CM curve $E:y^2=x^3-x$ itself, the GRH for $L(E,s)$ reduces to the GRH for Hecke $L$-functions of $\Q(i)$ via Deuring's factorization --- a problem that is structurally simpler than full RH and deeply connected to the arithmetic of the Gaussian integers.

Whether the self-dual bridge from lattice self-consistency to zero distribution can be made rigorous remains the central open question.

\begin{thebibliography}{20}
\bibitem{ftdspec}
W.\,J.\,cpaci-tani~III,
\emph{Foundational Ternary Dynamics: The Complete Specification}, v5.28, 2026.

\bibitem{watson1939}
G.\,N.\,Watson,
Three triple integrals,
\emph{Quart.\ J.\ Math.\ Oxford} \textbf{10} (1939), 266--276.

\bibitem{coateswiles}
J.\,Coates and A.\,Wiles,
On the conjecture of Birch and Swinnerton-Dyer,
\emph{Invent.\ Math.} \textbf{39} (1977), 223--251.

\bibitem{rubin}
K.\,Rubin,
The ``main conjectures'' of Iwasawa theory for imaginary quadratic fields,
\emph{Invent.\ Math.} \textbf{103} (1991), 25--68.

\bibitem{deuring}
M.\,Deuring,
Die Zetafunktion einer algebraischen Kurve vom Geschlechte Eins,
\emph{Nachr.\ Akad.\ Wiss.\ G\"ottingen} (1953), 85--94.

\bibitem{silverman}
J.\,H.\,Silverman,
\emph{The Arithmetic of Elliptic Curves},
Springer, 2nd ed., 2009.

\bibitem{titchmarsh}
E.\,C.\,Titchmarsh,
\emph{The Theory of the Riemann Zeta-function},
Oxford University Press, 2nd ed.\ (revised by D.\,R.\,Heath-Brown), 1986.

\bibitem{whittaker}
E.\,T.\,Whittaker and G.\,N.\,Watson,
\emph{A Course of Modern Analysis},
Cambridge University Press, 4th ed., 1927.

\bibitem{wiles}
A.\,Wiles,
Modular elliptic curves and Fermat's Last Theorem,
\emph{Ann.\ of Math.} \textbf{141} (1995), 443--551.
\end{thebibliography}
\end{document}
